\documentclass[11pt]{article}
\usepackage[T1]{fontenc}
\usepackage{lmodern,amsmath,amssymb,amsthm,mathtools,booktabs,microtype}
\usepackage[margin=1.01in]{geometry}
\usepackage[dvipsnames]{xcolor}
\usepackage[colorlinks,linkcolor=MidnightBlue,citecolor=MidnightBlue,urlcolor=MidnightBlue]{hyperref}
\usepackage{enumitem}
\usepackage{fancyhdr}
\pagestyle{fancy}\fancyhf{}\fancyhead[L]{\small Character identification from Dirichlet zeros}
\fancyhead[R]{\small Research note}\fancyfoot[C]{\thepage}
\setlength{\headheight}{14pt}\setlength{\emergencystretch}{2em}
\newtheorem{theorem}{Theorem}[section]
\newtheorem{lemma}[theorem]{Lemma}
\newtheorem{corollary}[theorem]{Corollary}
\newtheorem{proposition}[theorem]{Proposition}
\newcommand{\Z}{\mathbb Z}
\newcommand{\R}{\mathbb R}
\newcommand{\C}{\mathbb C}
\newcommand{\me}{\mathrm e}
\newcommand{\ii}{\mathrm i}
\newcommand{\dd}{\,\mathrm d}
\newcommand{\abs}[1]{\left|#1\right|}
\newcommand{\sinc}{\operatorname{sinc}}
\newcommand{\ord}{\operatorname{ord}}
\newcommand{\PP}{\mathcal P}
\newcommand{\GG}{\mathcal G}
\newcommand{\WW}{\mathcal W}
\title{Identifying a primitive Dirichlet character\\from a short window of zeros}
\author{Research note}
\date{27 November 2026}
\begin{document}
\maketitle
\begin{abstract}
For primitive Dirichlet characters of conductor at most $Q$, the signed multiset of zero ordinates up to height $C(\log Q)^2/\log\log Q$ determines the character. The constant is absolute, the order is unrestricted, and the primitive character of conductor one is included. We use the Dirichlet GRH theorem of the supplied Lean development. The argument detects a fixed angular difference at many primes of size $X\asymp(\log Q)^2$. The two-prime upper-bound sieve allows a local explicit-formula test to tolerate neighboring prime powers at resolution $(\log X)/X$ in the logarithmic prime variable. The resulting gain of two factors of $\log\log Q$ over the supplied quadratic separation bound persists under perturbations of the zero lists. We also give a translated-window theorem and compress the tests to $O(\log Q)$ complex measurements shared by all candidate characters.
\end{abstract}

\section{The separation theorem}
A zero ordinate is the signed number $\gamma$ in a nontrivial zero $1/2+\ii\gamma$. Write $\Gamma_\chi(I)$ for the multiset in the closed interval $I$, retaining every multiplicity. All logarithms are natural. For $Q\ge3$ put
\[
 L=\log(2Q),\qquad M=\log(2L),\qquad I(U,R)=[U-R,U+R].
\]
An occurrence at a repeated ordinate can be matched separately. For any partial one-to-one matching $\mathcal M$ between two finite multisets of real numbers, let $d_0$ be the total number of unmatched occurrences on the two sides. Define its cost at scale $M$ by
\[
 D_M(\mathcal M)=d_0+\sum_{(s,t)\in\mathcal M}\min\{2,M|s-t|\}.
\]
An exact common occurrence costs zero. Deleting and inserting one occurrence costs two in total.

\begin{theorem}[Uniform separation in a window]\label{th:main}
Use Dirichlet GRH. There are absolute constants $C,c>0$ such that, for $U\in\R$ and
\begin{equation}
 R\ge C\left(\frac{L^2}{M}+\log\log(\me^\me+|U|)\right),\label{eq:width}
\end{equation}
every two different primitive Dirichlet characters of conductors at most $Q$ satisfy
\begin{equation}
 D_M(\mathcal M)\ \ge\ c\,R\frac{M}{L}\label{eq:cost}
\end{equation}
for every partial matching of their multisets in $I(U,R)$.
\end{theorem}

At $U=0$ one can take $R=C' L^2/M$. The signed list then identifies the primitive character, its conductor and its order. Even the strictly positive interval $[R,3R]$ works by taking $U=2R$ and increasing the absolute constant. This latter observation uses the actual signs when selecting the zeros. Folding all ordinates about zero loses the distinction between a nonreal character and its conjugate.

\begin{corollary}[Multiplicity and displacement]\label{cor:mass}
In the window of Theorem~\ref{th:main}, write $m_j(t)$ for the multiplicity of $1/2+\ii t$ for character $j$. Then
\[
 \sum_{t\in I(U,R)} |m_1(t)-m_2(t)|\gg R M/L.
\]
If the lists have equal cardinalities, every bijection between them has $\sum|s-t|\gg R/L$. At the centered cutoff above, the first lower bound is $\gg\log(2Q)$ and the second is $\gg L/M$.
\end{corollary}
\begin{proof}
Match all exact common occurrences for the first assertion. For a bijection use $D_M\le M\sum|s-t|$.
\end{proof}
A shared zero with orders two and one contributes one unit to the first sum. Thus the question answered concerns multisets. An implication about a zero where precisely one function vanishes would require further information.

\paragraph{Sources and scope.}
The upstream Lean project \cite{lean} exports \texttt{LFunction\_eq\_zero\_re} for all Dirichlet characters, as well as \texttt{riemannHypothesis}. The former locates strip zeros on the critical line. We use it both for the two measured characters and for the primitive inducer of their quotient on the common group of units. The source obtains its half-plane theorem from a one-sided bound for a prime remainder. Here we start at the exported theorem and use a different explicit-formula argument. The project reports kernel checks with Lean 4.35.0 and mathlib pinned at \texttt{56a00585}. We inspected its statements and entry points. The downstream proofs in this note are on paper, and we did not rebuild the upstream project.

The ProveIt article \cite{proveit} merges 103 reports and explicitly describes its claims as unaudited by the merger. The section on recovery proves a quadratic separation height $O(L^2 M)$ by recovering every small prime-power coefficient, and asks whether the extra $\log\log Q$ can be removed and how to distinguish complex characters. Its section on direct consequences bounds the bias of small primes. We use the latter idea on a smooth dyadic interval and prove the estimates needed here. Its section on zeros gives stronger dependence on the central height for fixed characters at very large heights, with constants depending on the family. Our translated-window assertion keeps the conductor dependence uniform down to height zero.

\paragraph{Related work.}
Banks \cite{banks} distinguishes zero multisets unconditionally for same-modulus primitive characters in intervals of length $O_\theta(q^\theta\log T)$ above $T\gg_\theta q^\theta$, where $\theta>1/3$ (or $>1/4$ for cubefree $q$). Banks and Loftus \cite{bl} address genuinely unshared zeros at sufficiently large heights. Eid's preprint \cite{eid} already studies stable polylogarithmic-height determination for characters of orders dividing a prescribed integer using an isolated distinguishing coefficient. Its Theorem 1.1 gives $C_\varepsilon((\log(3Q))^2+d)^{1+\varepsilon}$ when both orders divide $d$. Priority for finite spectral identification or its bounded-order formulation is therefore not claimed here. The refinements pursued below are the conductor-uniform scale $L^2/M$, the removal of any order dependence, a linear-in-window separation bound for partially matched lists, and a common short list of measurements. We have not established an optimal height or priority against the entire literature.

\section{Many primes distinguish each pair}
Fix once and for all a smooth function $0\le w\le1$ of compact support in $(1,2)$ with $J=\int w(v)\dd v>0$. Let
\begin{equation}
 X=A L^2,\qquad \PP_X=\{p\text{ prime}: X<p\le2X\}, \qquad N=\#\PP_X,\label{eq:X}
\end{equation}
where the absolute constant $A$ will be fixed sufficiently large. Then $N\asymp X/\log X$ by the prime number theorem.

We record the precise explicit formula used in both arguments. For a sufficiently smooth compactly supported complex function $g$ on $(0,\infty)$, set
\[
 H_g(t)=\int g(u)\me^{\ii tu}\dd u,\qquad
 F_g(s)=\int g(u)\me^{(s-1/2)u}\dd u .
\]
For a primitive character $\chi$ with parity $\kappa\in\{0,1\}$, let $\delta_\chi=1$ for the character of conductor one, and zero otherwise. The positive-support explicit formula under Dirichlet GRH is
\begin{align}
 \sum_\gamma H_g(\gamma)
 &=-\sum_{n\ge2}\frac{\Lambda(n)\chi(n)}{\sqrt n}g(\log n)
   -\int g(u)\frac{\me^{-(\kappa+1/2)u}}{1-\me^{-2u}}\dd u
   +\delta_\chi\{F_g(0)+F_g(1)\}.\label{eq:EF}
\end{align}
Here $\Lambda(p^k)=\log p$ and $\Lambda(n)=0$ otherwise. One way to verify signs and normalization is to integrate $F_g(s)L'(s,\chi)/L(s,\chi)$ on a line $\Re s>1$. Mellin inversion gives the negative prime sum. Moving left crosses the nontrivial zeros and the trivial zeros at $-\kappa-2j$, whose transforms sum to the displayed integral. For zeta one adds back the spurious trivial zero at $0$ and includes the pole at $1$. Compact support away from $u=0$ gives exponential decay as the line is moved left, and integration by parts gives decay on horizontal segments chosen away from zeros. Equivalently this is the smooth Guinand--Weil formula with the negative prime channel and conductor term vanishing. The classical analytic ingredients, including the local zero count
\begin{equation}
 \#\{\gamma:|\gamma-t|\le1\}\ll\log(2q(|t|+3)),\label{eq:count}
\end{equation}
are standard \cite{mv,bmor,banks}. Formula \eqref{eq:EF} applies to the finitely differentiable windows below by smooth approximation, using \eqref{eq:count} to dominate the zero sum.

\begin{lemma}\label{lem:witness}
There is an absolute $b>0$ such that for every distinct pair $\chi_1,\chi_2$ in Theorem~\ref{th:main}, at least $bN$ primes in $\PP_X$ are coprime to both conductors and satisfy $|\chi_1(p)-\chi_2(p)|\ge1$.
\end{lemma}
\begin{proof}
For a primitive nonprincipal $\psi$ of conductor $q$, insert $g(u)=\me^{u/2}w(\me^u/X)$ in \eqref{eq:EF}. Then
$F_g(s)=X^s\int_1^2 w(v)v^{s-1}\dd v$. On the critical line the integral is $O_w((1+|\Im s|)^{-3})$, and the gamma integral is $O_w(1)$. Equation \eqref{eq:count} yields
\[
 \abs{\sum_n\Lambda(n)\psi(n)w(n/X)}\ll_w \sqrt X\log(2q).
\]
Prime powers of exponent at least two have total weight $O(\sqrt X\log^2 X)$ in $[X,2X]$, by counting their possible bases and exponents. Let $m$ be the least common multiple of the two conductors. The character $\chi_1\overline{\chi_2}$ induced to the units modulo $m$ is nonprincipal by uniqueness of a primitive inducer. Using its primitive inducer in the last bound and removing primes dividing $m$ costs at most $\log m$ after the higher powers have been removed. Here $m\le Q^2$. The prime number theorem also gives $\sum_{p\nmid m}\log p\,w(p/X)=JX+o(X)+O(\log m)$. Therefore
\[
 \sum_{p\nmid m} (\log p)w(p/X)|\chi_1(p)-\chi_2(p)|^2
 =2JX+O_w(\sqrt X\log(2Q))+o_w(X).
\]
Choose $A$ so that the right side exceeds $3JX/2$ and the total weight is at most $5JX/4$. This works for all $Q\ge3$ by increasing $A$. Each squared difference is at most four. The weight on differences of size at least one is consequently at least $JX/12$. Each prime carries weight at most $\log(2X)$, proving the claim.
\end{proof}

\section{Choosing prime probes at the mean spacing}
Fix $a=1/10$ and a large absolute $K$. Set $B_0=KX/\log X$. Define a quantity involving only ordinary prime powers by
\begin{equation}
 E(p)=\frac{\sqrt p}{\log p}
 \sum_{\substack{n\ne p,\,n\ge2\\ |\log(n/p)|\le a}}
 \frac{\Lambda(n)}{\sqrt n}
 \min\left(1,\frac{1}{B_0^2\log^2(n/p)}\right).
 \label{eq:leak}
\end{equation}
Its meaning is a worst-case error in measuring a single value $\chi(p)$ through all nearby coefficients. In particular, it can be computed before the character or the zero list is supplied.

\begin{lemma}\label{lem:sieve}
One can fix $K$, then increase $A$ if necessary, so that the set $\GG=\{p\in\PP_X:E(p)\le1/32\}$ omits at most $bN/2$ primes. Thus every pair in Lemma~\ref{lem:witness} differs by at least one on at least $bN/2$ primes of $\GG$.
\end{lemma}
\begin{proof}
On the indicated support $n\asymp p\asymp X$, each normalized von Mangoldt factor is bounded by an absolute constant. The mean value theorem gives $|\log(n/p)|\ge |n-p|/(3X)$. Put $D=3X/B_0$. The weight for an offset $d=|n-p|$ is bounded by $V_D(d)=\min(1,(D/d)^2)$.

For a prime neighbor with offset $d>0$, the pair upper-bound sieve \cite[Lemma 2.1]{evans} gives $O(X\mathfrak S(d)/\log^2X)$ ordered pairs in this dyadic range, with a constant covering both directions. Indeed the smaller prime lies between $\me^{-a}X$ and $2X$ and $d\le2(\me^a-1)X$. The odd shifts do not occur for large $X$. For even $d$, the singular series is a fixed constant times $\prod_{\ell\mid d,\ \ell>2}(1+1/(\ell-2))$. Its sum up to $Y$ is $O(Y)$, as follows by expanding the product and using $\prod_{\ell>2}(1+1/(\ell(\ell-2)))<\infty$. Splitting the offsets into $d\le D$ and dyadic shells above $D$ gives $\sum_d\mathfrak S(d)V_D(d)\ll D$ for $D\ge1$. We may ensure $D\ge1$ for all $Q$ by increasing $A$ after $K$ is chosen. Dividing by $N$, the mean contribution from primes is $O(D/\log X)=O(1/K)$.

There are $O(\sqrt X\log X)$ proper prime powers in $[\me^{-a}X,2\me^aX]$. For each such $n$, even summing over all integer centers $p\ne n$ gives $\sum_p V_D(|p-n|)\ll1+D$. Their mean contribution is $O((1+D)\log^2X/\sqrt X)=o(1)$. Choose $K$ and then $A$ to make the total mean less than $b/64$. Markov's inequality proves the assertion.
\end{proof}

The point of this lemma is the scale $1/B_0\asymp(\log X)/X$. A test whose compact support is forced to exclude every other integer has width only $O(1/X)$. At our wider resolution some neighboring primes really do contribute. Equation \eqref{eq:leak} bounds their combined effect uniformly in their character values. The sieve is used only for an upper bound.

\section{A localized explicit formula in any zero window}
Take $B=R/4\ge B_0$. Fix a small absolute $\lambda>0$, put $m=\lfloor\lambda B\rfloor$, and ensure $m\ge4$ by increasing the absolute constant in \eqref{eq:width}. With $\sinc v=\sin(v)/v$ and $\sinc 0=1$, set
\[
 \eta_m(v)=\begin{cases}\cos^{2m}(\pi v/(2a)),& |v|\le a,\\0,&|v|>a,\end{cases}
 \qquad \phi_B(v)=\eta_m(v)\sinc^2(Bv).
\]
For $p\in\GG$ take $y=\log p$ and
\begin{equation}
 g_p(u)=\me^{-\ii U(u-y)}\phi_B(u-y),\qquad
 H_p(t)=\me^{\ii t y}\int\phi_B(v)\me^{\ii(t-U)v}\dd v.\label{eq:probe}
\end{equation}
Thus $g_p(y)=1$, and moving the center $U$ leaves the absolute bounds on all prime weights unchanged. Given a zero list $S$ in $I=I(U,R)$ define
\begin{equation}
 A_p(S)=-\frac{\sqrt p}{\log p}\sum_{t\in S} H_p(t).\label{eq:read}
\end{equation}
All test functions are known from $Q,U,R,p$ alone.

\begin{lemma}[Truncation and coefficient error]\label{lem:read}
For every primitive $\chi$ of conductor at most $Q$ and every $p\in\GG$,
\[
 |A_p(\Gamma_\chi(I))-\chi(p)|\le 1/8 .
\]
Moreover $\|H_p\|_\infty\le\pi/B$ and $\|H_p'\|_\infty\le\pi(\log p+a)/B$ on $\R$.
\end{lemma}
\begin{proof}
The last bounds follow from $\int|g_p|\le\int_{\R}\sinc^2(Bv)\dd v=\pi/B$ and $|u|\le\log p+a$ on its support. In \eqref{eq:EF} the prime $p$ contributes precisely $\chi(p)$ after the normalization in \eqref{eq:read}. All other prime powers together cost at most $E(p)\le1/32$. The gamma kernel on the support is $O(X^{-1/2})$, so its normalized contribution is $O(1/(B\log X))$. The two pole terms for zeta cost $O(X/(B\log X))$. Taking $K$ sufficiently large makes the total archimedean and pole error at most $1/32$. Increasing $K$ in Lemma~\ref{lem:sieve} only improves its error bound.

It remains to control the unobserved zeros. The function $\sinc^2(Bv)$ is the characteristic function of the sum of two independent uniform random variables on $[-B,B]$. Hence the transform of $\phi_B$ averages translates, by at most $2B$, of the transform of $\eta_m$. On $[-a,a]$ the exponential expansion of $\cos^{2m}$ has nonnegative coefficients summing to one and frequencies of magnitude at most $\pi m/a$. The derivatives up to order $m-1$ vanish at the endpoints. Integrating by parts $m$ times gives, when $s=|t-U|\ge4B$,
\[
 |H_p(t)|\le2a\left(\frac{2\pi m}{as}\right)^m
      \ll \me^{-c_0B}(4B/s)^m.
\]
For instance choose $\lambda=a/(4\pi)$, which makes the base at $s=4B$ at most $1/8$. Applying \eqref{eq:count} on unit intervals in each dyadic shell $2^j4B\le s\le 2^{j+1}4B$ gives
\begin{equation}
 \sum_{|\gamma-U|>4B}|H_p(\gamma)|
 \ll B\me^{-c_0B}\log(2Q(3+|U|+B)).\label{eq:tail}
\end{equation}
Indeed the shell count is $O(2^jB[\log(2Q(3+|U|+B))+j])$ and the test decays by $2^{-jm}$. The implied constants are absolute, also for $q=1$.

For completeness, the growth conditions absorb every factor in \eqref{eq:tail}. From $B\ge K X/\log X$ we get $X\ll B^2$ and $\log(2Q)\ll B$. Taking $C$ in \eqref{eq:width} large gives $\log(3+|U|)\ll\exp(c_0 B/4)$. Thus \eqref{eq:tail}, even after multiplication by $\sqrt p/\log p$, is $O(B^3\exp(-c_0 B/2))$, which can be made at most $1/32$ uniformly. Adding the three error allowances proves the assertion.
\end{proof}

\begin{proof}[Proof of Theorem~\ref{th:main}]
Lemmas~\ref{lem:witness} and \ref{lem:sieve} give a $p\in\GG$ with $|\chi_1(p)-\chi_2(p)|\ge1$. Lemma~\ref{lem:read} forces the two values of $A_p$ to differ by at least $3/4$. In particular,
\[
 \left|\sum_{s\in\Gamma_1}H_p(s)-\sum_{t\in\Gamma_2}H_p(t)\right|
 \ge\tfrac34\,\frac{\log p}{\sqrt p},\qquad \Gamma_j=\Gamma_{\chi_j}(I).
\]
A partial matching bounds the left side by
\[
 \frac{\pi}{B}\left(d_0+\sum_{(s,t)\in\mathcal M}
       \min\{2,(\log p+a)|s-t|\}\right).
\]
Since $\log p+a\ll M$, $B=R/4$, and $(\log p)/\sqrt p\asymp M/L$, the two inequalities give \eqref{eq:cost}. All constants in the choice of $A,K,C$ are fixed independently of the characters, their orders, and the window.
\end{proof}

\paragraph{Error interpretation.}
If two possible true lists could both produce the same observed list by edits costing less than half of the lower bound in \eqref{eq:cost}, they would contradict the theorem. To see this, compose the two matchings through observed occurrences, keeping the pairs that survive both. Each unmatched endpoint costs at most the total deletion or insertion charges on its path. The truncated displacement $\min(2,M|x-y|)$ satisfies the triangle inequality. Thus the least edit cost for the true pair is at most the sum of the two costs. The lists being compared here always refer to the same prescribed window, including its endpoints.

\section{A common list of measurements}
The prime where two characters differ can depend on the pair. There is also a single short list that works simultaneously for the whole family. Each coordinate below is one complex weighted sum over the observed zero list.

\begin{theorem}[A common distinguishing list]\label{th:compression}
For each $Q$ there exists a list of $k\ll\log(2Q)$ primes $p_j\in\GG$, depending only on $Q$, and an absolute $\delta>0$, with the following properties.
\begin{enumerate}[label=(\roman*),leftmargin=2.1em]
\item Every two distinct primitive characters of conductor at most $Q$ satisfy $|\chi_1(p_j)-\chi_2(p_j)|\ge1$ on at least $\delta k$ coordinates.
\item For every window satisfying \eqref{eq:width}, the observations \eqref{eq:read} at these primes identify the unknown character. The same is true if fewer than $\delta k/2$ coordinates are arbitrarily corrupted and all remaining observations have additional absolute error at most $1/8$.
\end{enumerate}
\end{theorem}
\begin{proof}
Sample independently from the uniform distribution on $\GG$, allowing repeated indices. A fixed pair differs by at least one on a proportion at least $b/2$ of $\GG$. By the binomial lower-tail bound, the chance of fewer than $bk/4$ successful coordinates is at most $\exp(-bk/16)$. The number of characters is at most $\sum_{q\le Q}\varphi(q)\le Q^2$. A union bound over at most $Q^4$ pairs gives positive probability of simultaneous success for $k=\lceil C_1\log(2Q)\rceil$ with $C_1$ large. Put $\delta=b/4$.

For decoding, count for each candidate $\psi$ the number of coordinates whose observed complex value is farther than $1/4$ from $\psi(p_j)$. The true character has at most the number $e$ of corrupted coordinates counted against it, by Lemma~\ref{lem:read}. Every other character has at least $\delta k-e$ discrepancies, since at an uncorrupted separating coordinate the distance to the other candidate is at least $3/4$. The minimizing character is unique if $2e<\delta k$.
\end{proof}
The computation of a coordinate uses all zeros supplied in the window. The bound $O(\log Q)$ concerns the number of scalar measurements; it does not bound the number of individual zeros one must locate. The probabilistic proof gives existence of a common list and a sampling procedure. Precomputation can verify the separation requirement by testing candidate pairs, without any zero computations.

The logarithmic number of measurements has the right order in the bounded-amplitude, fixed-accuracy model. More generally, suppose each character is assigned $k$ complex coordinates of modulus at most two and that errors of size $1/8$ in every coordinate always permit unique decoding. Partition the square $[-2,2]^2$ into a fixed number $D$ of cells of diameter at most $1/8$. Two vectors with the same product cell can give the same noisy observation. Thus there are at most $D^k$ candidates. A prime $q\in(Q/2,Q]$ furnishes $q-2\gg Q$ primitive characters for large $Q$, whence $k\gg\log Q$. This is a statement about bounded scalar measurements with noise. The optimal zero height is a separate question.

\section{Numerical examples and remaining questions}
For the conjugate primitive quartic characters of conductor five, the nearest signed ordinates computed by the accompanying script are as follows. The underlying LMFDB records provide an independent comparison \cite{lmfdb}.
\begin{center}
\begin{tabular}{lrr}
\toprule
Character & Closest negative & Closest positive\\
\midrule
$\chi(2)=\ii$ & $-4.132903705213$ & $6.183578195451$\\
$\overline\chi(2)=-\ii$ & $-6.183578195451$ & $4.132903705213$\\
\bottomrule
\end{tabular}
\end{center}
The absolute-ordinate multisets agree identically by conjugation, even if one observes all zeros. In $[-5,5]$ the signed multisets already differ. The script refines sign changes of the real Hardy functions, checks the tabulated rounded ordinates and independently computes each character. The general proof assumes no spacing or simplicity condition.

A second experiment tests the prime side for the quartic conjugate pair modulo $101$. At $X=4096$, $B=16X/\log X$, there are 464 primes in $(X,2X]$. At 227 their character values differ by two. With a fixed smooth envelope supported in $|v|<0.1$, 202 of those 227 have actual normalized off-diagonal kernel mass at most $0.1$. At $p=7219$ the ideal difference is $2\ii$ and the measured prime sum is $1.98627419\ii$. The sum of the off-diagonal weights is $0.02181944$; twice this bounds the error in the difference. The finite experiment uses the actual weights, while \eqref{eq:leak} is a uniform upper envelope. Neither numerical example estimates the absolute constants in Theorem~\ref{th:main}.

\paragraph{Further research questions.}
How far below $L^2/M$ can the first separating ordinate lie uniformly? The present proof spends $L^2$ on finding a positive density of order-one coefficient differences, then exploits the mean prime spacing once. Cancellation among interfering prime terms could give finer spectral resolution. Can one obtain a comparably conductor-uniform bound for a genuinely unshared zero, or for both directions of multiplicity excess separately? A lower bound on the absolute difference of multiplicities alone leaves these questions open. Finally, computing small explicit values of $C$ and $c$ would turn the asymptotic theorem into a certified numerical cutoff at specified conductors.

\begin{thebibliography}{10}\small
\bibitem{lean} L. Sernau et al., \emph{GRH for Dirichlet $L$-functions, in Lean 4}, based on the work of L. Alp\"oge. Pinned \href{https://github.com/lsernau/grh-lean/tree/72fead102f663d5c7a407964a13dc0cf9d977207}{source tree}, \href{https://github.com/lsernau/grh-lean/blob/72fead102f663d5c7a407964a13dc0cf9d977207/grh/src/Main.lean}{exported theorems}, and \href{https://github.com/lsernau/grh-lean/blob/72fead102f663d5c7a407964a13dc0cf9d977207/grh/README.md}{verification report}.
\bibitem{proveit} V. Reshetnikov, \emph{ProveIt}, \href{https://github.com/VladimirReshetnikov/ProveIt/tree/f6b9918d87befd09f8786df70e668452171abb89/Analysis/GeneralizedRiemannHypothesis/Article}{merged GRH article}. See \href{https://github.com/VladimirReshetnikov/ProveIt/blob/f6b9918d87befd09f8786df70e668452171abb89/Analysis/GeneralizedRiemannHypothesis/Article/sections/rec.tex\#L2319-L2375}{quadratic coefficient recovery and separation}, \href{https://github.com/VladimirReshetnikov/ProveIt/blob/f6b9918d87befd09f8786df70e668452171abb89/Analysis/GeneralizedRiemannHypothesis/Article/sections/rec.tex\#L2617}{recovery questions}, and sections \texttt{dir.tex}, \texttt{zer.tex}.
\bibitem{eid} S. Eid, \emph{Effective Distinguishing of Primitive Dirichlet Characters by Their Low-Lying Zeros}, revised manuscript, July 2026, Theorems 1.1 and 1.3, \href{https://doi.org/10.5281/zenodo.21326653}{Zenodo}.
\bibitem{banks} W. D. Banks, \emph{Comparing zeros of distinct Dirichlet $L$-functions}, Publ. Math. Debrecen 106 (2025), 327--335. \href{https://doi.org/10.5486/PMD.2025.9899}{DOI}; \href{https://arxiv.org/abs/2302.07073}{arXiv}.
\bibitem{bl} W. Banks and K. Loftus, \emph{Unshared zeros of Dirichlet $L$-functions}, \href{https://arxiv.org/abs/2607.22930}{arXiv:2607.22930} (2026).
\bibitem{evans} N. Evans, \emph{Correlations of almost primes}, Math. Proc. Cambridge Philos. Soc. 174 (2023), Lemma 2.1. \href{https://doi.org/10.1017/S0305004122000251}{DOI}.
\bibitem{mv} H. L. Montgomery and R. C. Vaughan, \emph{Multiplicative Number Theory I. Classical Theory}, Cambridge University Press, 2007, Chapters 12--13.
\bibitem{bmor} M. A. Bennett, G. Martin, K. O'Bryant and A. Rechnitzer, \emph{Counting zeros of Dirichlet $L$-functions}, Math. Comp. 90 (2021), 1455--1482. \href{https://doi.org/10.1090/mcom/3599}{DOI}.
\bibitem{lmfdb} The LMFDB Collaboration, Dirichlet $L$-functions \href{https://www.lmfdb.org/L/1/5/5.2/r1/0/0}{\texttt{1-5-5.2-r1-0-0}} and \href{https://www.lmfdb.org/L/1/5/5.3/r1/0/0}{\texttt{1-5-5.3-r1-0-0}}.
\end{thebibliography}
\end{document}
